\subsection{Boundary conditions}
\label{sec:model-bc}

The system is closed by conditions at the two ends of the domain: the electrode surface, \(x=0\), and the edge of the near-electrode layer, \(x=L\), where the solution is matched to the adjacent LTE bulk plasma. The potential is referenced to zero at the electrode, \(\phi(0)=0\).

\emph{Electrode surface, \(x=0\).} Electrons reaching the surface are absorbed; they are assumed to have a Maxwellian distribution at their local temperature \(T_e\). At the cathode, thermionic emission of electrons, \(\Gamma_{\rm em}\), with a Maxwellian distribution at the wall temperature \(T_{\rm wall}\), is also included. The electron flux into the surface is
\begin{equation}
\boldsymbol{\Gamma}_e\cdot\hat{\mathbf{n}}=\tfrac14n_e\bar c_e-\Gamma_{\rm em},
\qquad
\bar c_e=\sqrt{\frac{8k_BT_e}{\pi m_e}} .
\label{eq:bc-electron}
\end{equation}
The emitted flux is given by the Richardson--Dushman law with Schottky lowering of the work function \(W\) by the field at the surface,
\begin{equation}
\Gamma_{\rm em}=\frac{A_RT_{\rm wall}^2}{e}\exp\!\left(-\frac{W-\Delta W}{k_BT_{\rm wall}}\right),
\qquad
\Delta W=\sqrt{\frac{e^3E_w}{4\pi\varepsilon_0}},
\qquad
E_w=\max(\mathbf{E}\cdot\hat{\mathbf{n}},0),
\label{eq:bc-emission}
\end{equation}
where \(A_R\) is the Richardson constant and \(W\) the work function. Each absorbed electron removes a mean energy \(2k_BT_e\) from the plasma, and each emitted electron brings in \(2k_BT_{\rm wall}\), so the boundary condition for the electron energy equation~(\ref{eq:eenergy}) is
\begin{equation}
\mathbf{h}_e\cdot\hat{\mathbf{n}}=2k_BT_e\,\tfrac14n_e\bar c_e-2k_BT_{\rm wall}\,\Gamma_{\rm em} .
\label{eq:bc-electron-energy}
\end{equation}

Ions reaching the surface are assumed to recombine there and to return to the plasma as atoms. The ion flux into the surface is taken as the sum of a thermal-effusion flux and a directed flux. Let \(\boldsymbol{\Gamma}_i^{d}\) denote the ion flux~(\ref{eq:flux-i}) without its density-diffusion term, \(\boldsymbol{\Gamma}_i^{d}=\boldsymbol{\Gamma}_i+[k_BT_h/(M\bar\nu^M_i)]\nabla n_i\). Ions are absorbed at their thermal-effusion flux plus this directed flux where it points into the wall,
\begin{equation}
\boldsymbol{\Gamma}_i\cdot\hat{\mathbf{n}}=
\tfrac14n_i\bar c_i+\max\!\left(\boldsymbol{\Gamma}_i^{d}\cdot\hat{\mathbf{n}},0\right),
\qquad
\bar c_i=\sqrt{\frac{8k_BT_{\rm wall}}{\pi M}} .
\label{eq:bc-ion}
\end{equation}
Where the sheath field drives the ions into the electrode, the directed term dominates. Elsewhere, the thermal-effusion flux ensures that ions reaching the wall are absorbed even where transport alone does not carry them into it.

The wall is impermeable to the heavy gas, \(\mathbf{u}\cdot\hat{\mathbf{n}}=0\); this is the boundary condition for the heavy-gas momentum equation~(\ref{eq:heavy-momentum}). Equation~(\ref{eq:flux-n}) then gives
\begin{equation}
\boldsymbol{\Gamma}_n\cdot\hat{\mathbf{n}}=-\boldsymbol{\Gamma}_i\cdot\hat{\mathbf{n}} :
\label{eq:bc-neutral}
\end{equation}
every ion neutralized at the surface re-enters the plasma as an atom. Secondary electron emission by ion impact is not included.

The heavy-particle temperature at the surface equals the electrode temperature, \(T_h(0)=T_{\rm wall}\). At the cathode, \(T_{\rm wall}\) is obtained by balancing the heat flux \(q_w\) delivered by the plasma against steady conduction through the cathode body,
\begin{equation}
q_w=\frac{k_w}{L_w}\left(T_{\rm wall}-T_0\right),
\label{eq:bc-heat-balance}
\end{equation}
where \(k_w\) is the thermal conductivity of the cathode, \(L_w\) its length, and \(T_0\) the temperature of its base. This crude heat balance is intended only to capture the trend of the cathode temperature with the current. The anode is usually water-cooled, so its surface temperature is assumed to be constant.

The heat flux into the surface is
\begin{equation}
q_w=\left(2k_BT_e+W\right)\tfrac14n_e\bar c_e
-\left(2k_BT_{\rm wall}+W\right)\Gamma_{\rm em}
+\left(\varepsilon_{\rm ion}-W\right)\boldsymbol{\Gamma}_i\cdot\hat{\mathbf{n}}
-\kappa_h\nabla T_h\cdot\hat{\mathbf{n}} .
\label{eq:bc-qwall}
\end{equation}
The first term is the energy delivered by the absorbed electrons: their kinetic energy plus the work function released on entering the metal.
 The second is emission cooling: each emitted electron removes the work function and its kinetic energy. The third is the energy released when an ion recombines at the surface, \(\varepsilon_{\rm ion}\), less the work function spent to extract the neutralizing electron. The last is heat conduction from the heavy gas. At the anode, \(\Gamma_{\rm em}=0\) and the emission term vanishes.

\emph{Bulk-plasma edge, \(x=L\).} The electron and ion densities and both temperatures are set to Dirichlet values from the adjacent LTE bulk plasma: \(n_e=n_e^{\rm bulk}\), \(n_i=n_i^{\rm bulk}\), and \(T_e=T_h=T_{\rm bulk}\); the last of these is the boundary condition for both energy equations. The neutral density follows from Dalton's law at the background pressure \(p_0\),
\begin{equation}
n_n=\frac{p_0-n_ek_BT_e}{k_BT_h}-n_i .
\label{eq:bc-neutral-density}
\end{equation}
This edge is open to the heavy gas: the heavy-gas momentum equation~(\ref{eq:heavy-momentum}) sees the bulk state, and hence the pressure \(p_0\), outside the boundary, and the heavy-gas velocity is given a zero normal gradient. The arc current is imposed through the electric field in Poisson's equation~(\ref{eq:poisson}): the total current density at the boundary, conduction plus displacement, is held equal to the applied current density \(J_{\rm app}\),
\begin{equation}
e\left(\boldsymbol{\Gamma}_i-\boldsymbol{\Gamma}_e\right)+\varepsilon_0\frac{\partial E}{\partial t}=J_{\rm app},
\label{eq:bc-current}
\end{equation}
which is integrated in time as an evolution equation for the boundary field. 